\section{问题三：未知类别玻璃的鉴别}

\subsection{未知样品预处理}
对表单3执行与训练数据完全一致的有效性检验、零值替代、闭合归一化、变量变换和标准化。所有预处理参数仅由已分类训练样本估计，再应用于A1--A8，避免使用未知样品信息反向影响模型。

\subsection{类别预测}
将未知样品输入问题二选定的分类器，输出预测类别、类别概率或判别距离。

\begin{table}[H]
  \centering
  \caption{未知样品类别鉴别结果}
  \label{tab:unknown-result}
  \small
  \begin{tabularx}{\textwidth}{C{1.2cm} C{2.2cm} C{2.2cm} Y C{2.0cm}}
    \toprule
    样品 & 风化状态 & 预测类型 & 属于铅钡的概率或判别得分 & 稳定性 \\
    \midrule
    A1 & \todo{填写} & \todo{填写} & \todo{填写} & \todo{填写} \\
    A2 & \todo{填写} & \todo{填写} & \todo{填写} & \todo{填写} \\
    A3 & \todo{填写} & \todo{填写} & \todo{填写} & \todo{填写} \\
    A4 & \todo{填写} & \todo{填写} & \todo{填写} & \todo{填写} \\
    A5 & \todo{填写} & \todo{填写} & \todo{填写} & \todo{填写} \\
    A6 & \todo{填写} & \todo{填写} & \todo{填写} & \todo{填写} \\
    A7 & \todo{填写} & \todo{填写} & \todo{填写} & \todo{填写} \\
    A8 & \todo{填写} & \todo{填写} & \todo{填写} & \todo{填写} \\
    \bottomrule
  \end{tabularx}
\end{table}

\subsection{分类敏感性分析}
对每个未知样品的成分进行多次小幅扰动，并在每次扰动后重新执行完整预处理与分类。定义样品\(i\)的稳定率为
\begin{equation}
  S_i=\frac{1}{B}\sum_{b=1}^{B}
  I\!\left(\hat{T}_i^{(b)}=\hat{T}_i^{(0)}\right),
\end{equation}
其中\(B\)为扰动次数。若\(S_i\)较低或预测概率接近阈值，则将该样品标记为边界样品，并重点讨论其不确定性。

同时比较逻辑回归、LDA、决策树或SVM等候选模型对A1--A8的预测一致性。\todo{列出发生分歧的样品及原因。}
